\chapter{Architecture}
\label{ch:architecture}

\section{Shape of the system}

Big\_Agent is a single FastAPI process over a single SQLite file. There is no
message broker, no worker fleet and no external index: background work runs in
an in-process thread pool, and the vector index is a NumPy matrix rebuilt from
the database when it goes stale. That is a deliberate constraint --- the
application has to be installable on one engineer's workstation by copying a
folder.

\begin{shellblock}
HTTP client (SmartCAE v2 UI, legacy UI, scripts, other machines on the LAN)
     |
     v
+---------------------------------------------------------------+
| app/main.py -- FastAPI: routes, auth middleware, OpenAPI       |
+---------------------------------------------------------------+
     |                   |                      |
     v                   v                      v
+-----------+    +----------------+    +----------------------+
| parsers   |    | services       |    | job_manager          |
| pdf/docx/ |    | ingest, search |    | background jobs      |
| pptx      |    | qa, review ... |    | (thread pool)        |
+-----------+    +----------------+    +----------------------+
     |                   |                      |
     v                   v                      |
+-----------+    +----------------+             |
| processing|    | vector_index   |             |
| clean,    |    | NumPy matrix   |             |
| chunk     |    | + cache stamp  |             |
+-----------+    +----------------+             |
     |                   |                      |
     +---------+---------+----------------------+
               v
     +--------------------------+
     | db: SQLAlchemy + SQLite  |
     | data/app.db (WAL)        |
     +--------------------------+
               |
               v
     data/documents/  -- stored report files, named by hash
\end{shellblock}

\section{Package map}

\begin{longtable}{L{0.28\textwidth} L{0.62\textwidth}}
\toprule
\textbf{Module} & \textbf{Responsibility} \\
\midrule
\endfirsthead
\toprule
\textbf{Module} & \textbf{Responsibility} \\
\midrule
\endhead
\bottomrule
\endfoot
\file{app/main.py} & Every HTTP route, the auth middleware, the OpenAPI
customization and the chat-routing heuristics. The only module that knows about
requests. \\
\file{app/config.py} & \code{Settings} --- the single source of configuration.
See Chapter~\ref{ch:configuration}. \\
\file{app/branding.py} & \code{AppBrand} records: display name, initials,
palette, cookie and data directory defaults per product variant. \\
\file{app/version.py} & \code{APP\_VERSION}, reported by \route{/health} and used
as the FastAPI application version. \\
\file{app/schemas.py} & The internal pipeline dataclasses:
\code{ParsedSection}, \code{CleanSection}, \code{ChunkPayload}. Not HTTP
models. \\
\file{app/api\_models.py} & The pydantic request and response models for the
HTTP layer. \\
\file{app/db/} & \file{models.py} (ORM tables) and \file{session.py} (engine,
pragmas, additive schema evolution, the \code{get\_session} dependency). \\
\file{app/parsers/} & One function per input format, all returning
\code{list[ParsedSection]}. \\
\file{app/processing/} & \code{text\_cleaner} (normalize, drop repeated
headers/footers), \code{chunker} (overlapping word blocks),
\code{extraction\_metrics} (per-page quality rollup). \\
\file{app/text/normalize.py} & Turkish-aware folding and tokenization shared by
search, catalog matching and review rules. \\
\file{app/rules/} & \file{profile\_catalog.py} plus one JSON file per discipline
under \file{profiles/}. Discipline rules are data, not code. \\
\file{app/services/} & Everything else: ingest, search, retrieval, QA, catalog,
review, comparison, drafting, duplicates, graph, jobs, the CATIA bridge. \\
\file{app/ui/}, \file{app/static/} & Front ends: SmartCAE v2 workspace, the two
landing pages, and the legacy single-page UI. \\
\end{longtable}

\section{Request lifecycle}

\begin{enumerate}
\item \textbf{Lifespan.} \code{init\_db()} runs once at startup:
\code{create\_all} plus the additive column check described below.

\item \textbf{Auth middleware.} When \env{APP\_AUTH\_ENABLED} is on, every
request except \route{/health}, \route{/login}, \route{/logout} and
\route{/favicon.ico} --- plus the landing page for the \code{raporhub} and
\code{repocto} variants --- must carry a valid signed session cookie. An HTML
\code{GET} without one is redirected to \route{/login}; anything else gets
\code{401}.

\item \textbf{Route.} Routes are ordinary \code{def} functions, so Starlette runs
them in its thread pool. That is why the SQLite engine is built with
\code{check\_same\_thread=False} and a 30-second busy timeout.

\item \textbf{Session dependency.} \code{get\_session} yields a
\code{SessionLocal} and closes it in a \code{finally} block. Services take the
session as a constructor argument --- they never open their own, except inside
background jobs.

\item \textbf{Response model.} Almost every route declares a response model from
\file{app/api\_models.py}, which is what makes the generated OpenAPI schema
worth reading.
\end{enumerate}

\section{Database engine}

SQLite is used with settings chosen for a multi-threaded single process.

\begin{longtable}{L{0.32\textwidth} L{0.58\textwidth}}
\toprule
\textbf{Setting} & \textbf{Why} \\
\midrule
\endfirsthead
\bottomrule
\endfoot
\code{check\_same\_thread=False} & Routes run in a thread pool, so connections
cross threads. \\
\code{timeout=30} / \code{PRAGMA busy\_timeout=30000} & A writer waits out a
concurrent write instead of failing with ``database is locked'' at the 5-second
default. \\
\code{PRAGMA journal\_mode=WAL} & Readers keep working during writes ---
important while a batch ingest job is running. \\
\code{PRAGMA synchronous=NORMAL} & Throughput, accepting the usual WAL
durability trade-off. \\
\code{autoflush=False} & Services decide when to flush; ingest depends on
explicit flush points to obtain generated ids. \\
\end{longtable}

\section{Schema evolution without a migration tool}
\label{sec:schema-evolution}

The project carries no migration framework, and \code{create\_all} only creates
missing \emph{tables}, never missing columns. A workstation whose
\file{data/app.db} predates a column would fail every query that selects it.

\file{app/db/session.py} therefore keeps an explicit map of columns added to
tables that already shipped, and \code{init\_db()} adds any that are absent with
\code{ALTER TABLE}:

\begin{pyblock}
_SQLITE_ADDED_COLUMNS = {
    "documents": {"extraction_quality": "JSON"},
    "document_pages": {
        "extraction_method": "VARCHAR(32)",
        "ocr_attempted": "BOOLEAN",
        "char_count": "INTEGER",
        "word_count": "INTEGER",
    },
}
\end{pyblock}

Two rules make this safe, and both must hold for any future entry:

\begin{itemize}
\item every added column stays \textbf{nullable and default-free} --- the rows
already on disk have no value to backfill; and
\item \textbf{readers treat NULL as ``unknown''}, never as a synthesized value.
A page with \code{extraction\_method IS NULL} was ingested before provenance
existed; it is not ``native with empty text''.
\end{itemize}

\section{Background jobs}

Long operations return immediately instead of holding an HTTP connection open.
\code{JobManager} (\file{app/services/job\_manager.py}) is a small in-process
scheduler over a \code{ThreadPoolExecutor} with a \textbf{single worker}, so jobs
run one at a time and never contend for the SQLite write lock.

\begin{shellblock}
POST /ingest/batch        -> 202 {"job_id": ..., "status": "queued"}
POST /catalog/ingest-sample
POST /catalog/ingest-selected
POST /duplicates/scan
POST /embeddings/rebuild

GET  /jobs/{job_id}       -> queued | running | succeeded | failed
                             + progress {done, total, message}
                             + result (on success) or error (on failure)
GET  /jobs                -> the most recent records
\end{shellblock}

A job body receives a \code{JobContext} and calls
\code{context.set\_progress(done, total, message)} as it advances. Jobs open
their \textbf{own} session (\code{SessionLocal()}) and close it themselves: the
request-scoped session is gone by the time the job runs. Finished records are
pruned so the list stays bounded.

\section{Vector index}

Semantic search does not query embeddings row by row.
\file{app/services/vector\_index.py} loads all stored vectors into one NumPy
matrix, normalizes it once, and answers a query with a single matrix product.

The cache is stamped with a cheap signature over \code{chunk\_embeddings}: the
row count, the maximum chunk id and the sum of chunk ids. That triple changes
whenever an embedding is added, removed or replaced, so a reader is never served
a stale matrix even if nobody invalidated it. Writers still call
\code{invalidate\_vector\_index()} after a commit --- ingest and the embedding
rebuild both do --- to drop the matrix eagerly.

Embeddings are stored as packed little-endian \code{float32} bytes. Databases
written by older versions hold JSON text in the same column;
\code{EmbeddingService.deserialize} accepts both, and one
\route{/embeddings/rebuild} run converts them.

\section{Service composition}

Services are plain classes constructed per request with an explicit session, and
they compose by construction rather than by inheritance.

\begin{shellblock}
DocumentIntelligenceService     -- chat/report answering
  +-- SearchService             -- keyword / semantic / hybrid
  |     +-- EmbeddingService    -- token-hash | sentence-transformers
  |     +-- VectorIndex         -- cached NumPy matrix
  +-- HaystackRetrievalService  -- the v3 retrieval implementation
  +-- ReportReviewService       -- rule findings, when the question asks
  +-- ReportQualityService      -- caption/numbering questions
  +-- LLMProvider               -- optional generation

RetrievalOrchestrator           -- opt-in enhancement shell
  +-- QueryUnderstandingService -- intent, expansions, metadata filters
  +-- SearchService
  +-- Reranker                  -- no-op unless enabled
\end{shellblock}

Every optional dependency has a null implementation (\code{NoOpReranker}, a
disabled \code{LLMProvider}, the \code{token-hash} embedding service), which is
how the application keeps working with nothing installed.
